% Generated from locale/tr/content/proof-theory/normalization/introduction.tex; do not hand-edit.


\olfileid[tr]{pt}{nor}{int}

\olsection{Giriş}

Bir koşullu~$!A \lif !B$ !!{prove}{}dıysanız, onu~\Elim{\lif} kullanarak
$!B$'nin !!a{proof}{}ında kullanabilirsiniz. Bunun için elbette $!A$'nın bir
!!a{proof}{}ı~$\delta_1$ da gerekir. $!A \lif !B$'nin !!{proof}{}ı~$\delta$,
büyük olasılıkla, açık $!A$ varsayımından $!B$'nin bir !!{proof}{}ı~$\delta_2$
üretip onu $!A \lif !B$'nin ispatına dönüştüren~\Intro{\lif} ile biter. Öyle
ise son !!{proof}{}ınız şu biçimdedir:
\[
\AxiomC{$\Discharge{!A}{x}$}
\RightLabel{$\delta_2$}
\DeduceC{$!B$}
\DischargeRule{\Intro{\lif}}{x}
\UnaryInfC{$!A \lif !B$}
\AxiomC{}
\RightLabel{$\delta_1$}
\DeduceC{$!A$}
\RightLabel{\Elim{\lif}}
\BinaryInfC{$!B$}
\DisplayProof
\]
Stratejiniz elinizde bulunan bir şeyi ($!A \lif !B$'nin !!a{proof}{}ını)
yeniden kullandığı için verimli olsa da !!{proof}{}ınız verimli değildir.
$!A \lif !B$'yi yalnızca hemen gidermek üzere tanıtmak bir dolambaç gibi
görünür. $!B$'yi !!{prove}{}mak için daha doğrudan bir yol olmalıdır.

Gerçekte her zaman vardır. Yukarıdaki gibi ``dolambaçlar''---hemen ardından
!!a{elimination} kuralı gelen !!a{introduction} kuralı---doğal tümdengelim
!!{proof}s{}ından her zaman çıkarılabilir. Buna \emph{normalleştirme}, elde
edilen ispata da \emph{normal} !!{proof} (ya da \emph{normal biçimde}
!!a{proof}) denir.

Bu sonucun ispatı, sekant hesabındaki kesme-giderimi teoremi gibi biraz
karmaşıktır ve çok sayıda durumun ele alınmasını gerektirir. Klasik
$\Log{N1c}$ için ispatı, sezgici $\Log{N1i}$ için ispatından daha zordur.
Kesme-giderimi, sekant hesabındaki \Cut{} kuralıyla daha fazla şey
!!{prove}{}yamayacağınızı gösterdiği gibi normalleştirme de dolambaçlar
kullanarak, onlardan kaçınarak ispatlanabilenden daha fazla şey
ispatlayamayacağınızı gösterir. Başka benzerlikler de vardır: Normal bir
!!{proof}, aynı sonucun dolambaçlı en kısa !!{proof}{}ından çok daha uzun
olabilir; tıpkı sekant hesabındaki kesmeli !!{proof}s{}ın en kısa kesmesiz
!!{proof}{}tan çok daha kısa olabilmesi gibi. Bir sonucu !!{proof}{}ınızda
yalnızca sonucun alt-!!{formula}s{}ini kullanarak her zaman !!{prove}{}yabileceğinizi
söyleyen alt-!!{formula} özelliği, sekant hesabında
kesme-gideriminden çıktığı gibi doğal tümdengelimde normalleştirmeden çıkar.

